\section{A quotient obtained by cutting off a band}
\label{sec:band-quotient}

The next step transfers support from a smaller cyclic staircase module
and a band into $W_m$. The projection must retain their specified
generators and realize the full induced module, not just a submodule
of it. We first define the two factors and then verify both requirements.

Fix $s\in\{1,2\}$ and $h=m-s\ge1$, and put
\begin{equation}\label{band:sets}
 D=B_h,\qquad F=B_m\setminus D,\qquad
 G=G_m=S_{B_m},\qquad H=S_D\times S_F.
\end{equation}
Thus $|F|=N_m-N_h=s(2m-s+1)/2$.
Let $R_F$ and $C_F$ be the partitions of $F$ into its nonempty row
and column intersections, with the orders inherited from $R_m$ and $C_m$.
Each block has size at most $s$.  In two layers with alphabet
$1,\ldots,s$, define
\begin{equation}\label{band:generator}
 u_{m,s}=v_{R_F}\otimes v_{C_F},\qquad
 U_{m,s}=\C[S_F]u_{m,s}.
\end{equation}
The dependence on $m$ records the length of the band.
Our goal is a quotient of $W_m$ induced from $W_h\boxtimes U_{m,s}$.

In the disjoint-coset model of induction
\cite[Section~3.3]{fultonharris1991}, translates of the inducing module
form a direct sum. We record why disjoint coordinate sectors realize
this model for the cyclic span of a specified vector, including the
inverse map.

\begin{lemma}[Induction from coordinate sectors]\label{lem:sector-induction}
Let a finite group $G$ act transitively on a finite set $\Omega$.
Suppose a $G$-module has a direct-sum decomposition
$T=\bigoplus_{A\in\Omega}T_A$ with $gT_A=T_{gA}$.
Fix $D\in\Omega$, let $H$ be its stabilizer, and take $z\in T_D$.
If $A_0=\C[H]z$, then the map
\begin{equation}\label{band:induced-map}
 \Phi:\C[G]\otimes_{\C[H]}A_0\longrightarrow\C[G]z,
 \qquad g\otimes a\longmapsto ga
\end{equation}
is an isomorphism of $G$-modules.
\end{lemma}

\begin{proof}
The relation $(gh)\otimes a=g\otimes ha$, for $h\in H$, has the
same image on both sides, so $\Phi$ is well-defined.  It is equivariant
and surjective.  Choose $g_A\in G$ with $g_AD=A$ for every $A\in\Omega$,
and let $\pi_A:T\to T_A$ be the coordinate projection.  Since
\[
 \C[G]z=\bigoplus_{A\in\Omega}g_AA_0,
 \qquad g_AA_0\subseteq T_A,
\]
we have $g_A^{-1}\pi_A(y)\in A_0$ for $y\in\C[G]z$.  Define
\begin{equation}\label{band:induced-inverse}
 \Psi(y)=\sum_{A\in\Omega}
          g_A\otimes g_A^{-1}\pi_A(y).
\end{equation}
Replacing $g_A$ by $g_Ah$ does not change this expression, because
$g_Ah\otimes h^{-1}a=g_A\otimes a$ in the balanced tensor product.
Clearly $\Phi\Psi(y)=\sum_A\pi_A(y)=y$.
For $g\otimes a$, set $A=gD$ and write $g=g_Ah$ with $h\in H$.
Only its $A$-sector is nonzero, and hence
\[
 \Psi\Phi(g\otimes a)=g_A\otimes ha=g\otimes a.
\]
Thus the displayed maps are inverse isomorphisms.
\end{proof}

\begin{proposition}[Band quotient]\label{prop:band-quotient}
For the parameters in \Cref{band:sets}, there is a surjective
$G_m$-equivariant map
\begin{equation}\label{band:quotient}
 W_m\twoheadrightarrow
 \Ind_{S_D\times S_F}^{S_{B_m}}(W_h\boxtimes U_{m,s}).
\end{equation}
In particular, every irreducible constituent of the induced module on
the right is a constituent of $W_m$.
\end{proposition}

\begin{proof}
\emph{Projection and the high positions.}
Use the ordinary tensor identification
\[
 E_m^{\otimes B_m}\otimes E_m^{\otimes B_m}
 \cong (E_m\otimes E_m)^{\otimes B_m}.
\]
Put $E_{\mathrm{lo}}=\Span(e_1,\ldots,e_s)$ and
$E_{\mathrm{hi}}=\Span(e_{s+1},\ldots,e_m)$.
On a pair-letter basis vector define
\[
 p_s(e_a\otimes e_b)=
 \begin{cases}
 e_a\otimes e_b,& a,b\le s\text{ or }a,b>s,\\
 0,&\text{otherwise},
 \end{cases}
 \qquad P_s=p_s^{\otimes B_m}.
\]
The same local projection is used at every position, so $P_s$
commutes with $G$.  Call a retained position \emph{high} when both
letters are greater than $s$.

In a retained word of $w_m$, let $A$ be the set of high positions.
Alternation in the first layer gives row counts
$r_i=\max(h+1-i,0)$, and alternation in the second gives column
counts $d_j=\max(h+1-j,0)$, for $1\le i,j\le m$.
For each $1\le k\le m$, these counts satisfy
\begin{equation}\label{band:margin-equality}
 \sum_{i=1}^k r_i
 =\#\{(i,j)\in B_h:i\le k\}
 =\sum_{j=1}^m\min(k,d_j).
\end{equation}
Column $j$ can contribute at most $\min(k,d_j)$ high positions to
the first $k$ rows.  Equality of the sums in
\Cref{band:margin-equality} therefore forces equality in every column.
For $d_j>0$, take $k=d_j$: all its $d_j$ high positions occupy the
first $d_j$ rows.  Columns with $d_j=0$ have no high positions.
Consequently
\begin{equation}\label{band:unique-high-set}
 A=\{(i,j)\in B_m:i\le h+1-j\}=B_h=D.
\end{equation}
For the zero-one matrix recording membership in $A$, this is the
equality case of the row and column margin inequalities
\cite[Section~2]{ryser1957}. The specific staircase margins in
\Cref{band:margin-equality} make every prefix bound an equality.

\emph{Factorization with signs.}
In a row or column block of length $r$, its intersection with $D$
is its initial segment of length $a=\max(r-s,0)$; its intersection
with $F$ has length $b=\min(r,s)$.
The retained assignments put the high letters $s+1,\ldots,r$ on
the first segment and the low letters $1,\ldots,b$ on the second.
Let $\Alt$ denote the unnormalized alternation of a displayed list,
placed in the indicated segment, with $\Alt(\varnothing)=1$.
The part of the block tensor with this prescribed high set is exactly
\begin{equation}\label{band:block-factorization}
 (-1)^{ab}\,
 \Alt(e_{s+1},\ldots,e_r)\otimes\Alt(e_1,\ldots,e_b).
\end{equation}
Here the high list is empty when $r\le s$.
Indeed, every high letter precedes every low letter in the position
order, giving $ab$ inversions; all other inversions are internal to
the two independently permuted lists.  Every such pair of permutations
occurs once.  In particular, when $r<s$ the formula is just the
original alternation on $e_1,\ldots,e_r$.

Rows and columns both have the length list $m,m-1,\ldots,1$.
Thus their cross signs in \Cref{band:block-factorization} multiply to
\[
 \prod_{r=1}^m(-1)^{2\max(r-s,0)\min(r,s)}=1.
\]
Identify $E_h$ with $E_{\mathrm{hi}}$ by
$e_a\mapsto e_{s+a}$, and write $\widetilde w_h$ and
$\widetilde W_h$ for the images of $w_h$ and $W_h$ in the two
high-letter layers on $D$.  The low portions are precisely the blocks
of $R_F,C_F$.  Regrouping the factors over the disjoint sets $D,F$
introduces no signs, and hence
\begin{equation}\label{band:generator-factorization}
 z:=P_s(w_m)=\widetilde w_h\otimes u_{m,s}\ne0.
\end{equation}
The nonvanishing follows also directly from the nonzero alternating
tensors on each disjoint block.

The two factors of $H=S_D\times S_F$ act independently on this
tensor, while each acts diagonally on the two layers over its own
position set.  Therefore
\begin{align}
 A_0:=\C[H]z
 &=\Span\{(\sigma\widetilde w_h)\otimes(\tau u_{m,s}):
             \sigma\in S_D,\ \tau\in S_F\}\notag\\
 &=\widetilde W_h\boxtimes U_{m,s}
 \cong W_h\boxtimes U_{m,s}.
 \label{band:subgroup-orbit}
\end{align}
The second equality follows by expanding tensors of linear
combinations of the two separate orbit sets.

\emph{The coordinate sectors.}
Let $\Omega$ be the set of subsets $A\subset B_m$ of size $|D|$.
Within the pair-letter space define
\[
 T_A=(E_{\mathrm{hi}}\otimes E_{\mathrm{hi}})^{\otimes A}
       \otimes
       (E_{\mathrm{lo}}\otimes E_{\mathrm{lo}})^{\otimes(B_m\setminus A)}.
\]
These are the coordinate spans of words high exactly on $A$.
They have disjoint word bases, so their sum is direct, and
$gT_A=T_{gA}$.  By \Cref{band:generator-factorization}, $z\in T_D$;
the stabilizer of $D$ in $G$ is exactly $H$.
Apply \Cref{lem:sector-induction} and
\Cref{band:subgroup-orbit} to obtain
\[
 \C[G]z\cong\Ind_H^G A_0
       \cong\Ind_H^G(W_h\boxtimes U_{m,s}).
\]
Since $P_s$ is equivariant, $P_s(W_m)=\C[G]P_s(w_m)=\C[G]z$.
Its restriction to $W_m$, followed by the inverse map $\Psi$ in
\Cref{band:induced-inverse} and the identification of $A_0$ in
\Cref{band:subgroup-orbit}, gives \Cref{band:quotient}.
The image of this ambient projection need not lie in the two original
Specht factors; the composition is still a surjective $G$-equivariant
linear map onto
the displayed module.  Semisimplicity over $\C$ proves the last
assertion.
\end{proof}
